\subsection{Equações de Hill/CW}

Como \eqref{eq:relativo_LVLHgeral} estava sendo expresso no referencial do alvo, vetor posição é:
\begin{equation}
\label{eq:posicao_LVLH_target}
\vec r_t
=
\begin{bmatrix}
0 \\
0 \\
-r
\end{bmatrix},
\end{equation}
seguindo a orientação dextrogiro.

E também o vetor de velocidade angular:
\begin{equation}
\label{eq:velocidadeAngular_LVLH_target}
\vec\omega=
\begin{bmatrix}
0 \\
-\omega \\
0
\end{bmatrix}
\end{equation}

Aplicando os produtos vetoriais apresentados em \eqref{eq:relativo_LVLHgeral}, obtêm-se:

\begin{equation}
\label{eq:velAngular_cross_s}
\vec\omega \times \vec s^{\,*} =
\begin{bmatrix}
-\omega z \\
0         \\
\omega x
\end{bmatrix} 
\end{equation}

\begin{equation}
\label{eq:velAngular_cross_cross_s}
\vec\omega \times (\vec\omega \times \vec s^{\,*} )
=
\begin{bmatrix}
-\omega^2 x \\
0           \\
-\omega^2 z 
\end{bmatrix}
\end{equation}

\begin{equation}
\label{eq:omega_cross_dsdt}
\vec\omega \times \dot{\vec s}^{\,*} 
=
\begin{bmatrix}
-\omega \dot z \\
0 \\
\omega \dot x
\end{bmatrix}
\end{equation}

\begin{equation}
\label{eq:dOmega_cross_s}
\frac{d\vec\omega}{dt} \times \vec s^{\,*} 
=
\begin{bmatrix}
-\dot\omega z \\
0              \\
\dot\omega x
\end{bmatrix}
\end{equation}

\begin{equation}
\label{eq:jacobiana_LVLH}
\vec{M}s^{\,*} 
=
\begin{bmatrix}
    1 & 0 & 0 \\
    0 & 1 & 0  \\
    0 & 0 & -2
\end{bmatrix}
\vec s^{\,*} 
=
\begin{bmatrix}
x \\
y \\
-2z
\end{bmatrix}
\end{equation}

% No caso da órbita circular ($e=0$), a velocidade angular $\omega$ é constante e coincide com o movimento médio (\emph{mean motion}, $\eta$), definido por \cite{fehse2003}:

Como o momento angular $\vec{L}$ é constante para órbitas fixas elípticas e para o caso especial de órbitas circulares, a velocidade angular também é constante, sendo expressa por 

\begin{equation}
\label{eq:omega_constante1}
\omega^2 = \frac{\mu}{r_t^3}
\end{equation}

Como $r_t=a$, a velocidade angular $\omega^{2}= \eta^2 = \mu/r_t^{3}$, portanto, neste caso em particular de órbita circular, $\omega=n$.
Portanto, $\dot{\omega}=0$, pode ser reescrita como

\begin{equation}
\label{eq:omega_constante2}
\eta^2 = \frac{\mu}{r_t^3}
\end{equation}

Inserindo este resultado em \eqref{eq:relativo_LVLHgeral}, obtêm-se as equações linearizadas do movimento relativo, conhecidas como equações de Hill/Clohessy--Wiltshire \cite{hill1878,fehse2003,wie2008space}:

\begin{equation}
    \label{eq:equacoes_HillX}
\ddot{x} - 2\omega \dot{z} = \frac{F_x}{m_c}
\end{equation}

\begin{equation}
    \label{eq:equacoes_HillY}
\ddot{y} + \omega^2 y = \frac{F_y}{m_c}
\end{equation}

\begin{equation}
    \label{eq:equacoes_HillZ}
\ddot{z} + 2\omega \dot{x} - 3 \omega^2 z = \frac{F_z}{m_c}
\end{equation}

As equações em \eqref{eq:equacoes_HillX} a \eqref{eq:equacoes_HillZ} evidenciam os termos de Coriolis ($\pm2\omega$ sobre as velocidades cruzadas) e o termo centrípeto ou gravidade--gradiente (multiplicadores de $\omega^2$ sobre $x$ e $z$).
Ainda, essas equações diferenciais lineares com coeficientes que não variam no tempo representam o modelo geral que é válido para uma trajetória relativa arbitrária \cite{debruijn2011,ortolano2017} entre o perseguidor e o alvo, considerando que este último se move sob a influência exclusiva de um campo gravitacional central, ou seja, pode ser considerado uma \emph{espaçonave não-cooperativa} \cite{mahendrakar2023}.

Para reduzir a ordem das matrizes, as equações de Hill foram reescritas na forma de espaço de estados, expressa como \cite{fehse2003,starek2016,sherrill2014}

\begin{equation}
    \dot{\vec{x}} = \vec{A}\,\vec{x} + \vec{B}\,\vec{u}
    \label{eq:state_space_general}
\end{equation}

Onde $\vec{A}$ é a matriz de transição, $\vec{x}$ é o vetor de estados, $\vec{B}$ é a matriz de entradas e $\vec{u}$ é o vetor de entradas, que neste caso é a força de controle aplicada no perseguidor.
Para o primeiro caso dentro do plano (\textit{in-plane}), as equações \eqref{eq:equacoes_HillX} a \eqref{eq:equacoes_HillZ}, e considerando o vetor de estados
$[x, z, \dot{x}, \dot{z}]^T$ podem ser escritas conforme:

\begin{equation}
\begin{bmatrix}
\dot{x}(t) \\[0.3em]
\dot{z}(t) \\[0.3em]
\ddot{x}(t) \\[0.3em]
\ddot{z}(t)
\end{bmatrix}
=
\begin{bmatrix}
0 & 0 & 1 & 0 \\
0 & 0 & 0 & 1 \\
0 & 0 & 0 & 2\omega \\
0 & 3\omega^2 & -2\omega & 0
\end{bmatrix}
\begin{bmatrix}
x(t) \\[0.3em]
z(t) \\[0.3em]
\dot{x}(t) \\[0.3em]
\dot{z}(t)
\end{bmatrix}
+
\begin{bmatrix}
0 & 0 \\
0 & 0 \\
\frac{1}{m_c} & 0 \\
0 & \frac{1}{m_c}
\end{bmatrix}
\begin{bmatrix}
F_x \\[0.3em]
F_z
\end{bmatrix}
\label{eq:inplane}
\end{equation}

De forma similar, para a dinâmica fora do plano orbital, (\textit{out-of-plane}), considerando o vetor de estados$[y,\dot{y}]^T$, obtém-se:

\begin{equation}
\begin{bmatrix}
\dot{y}(t) \\[4pt]
\ddot{y}(t)
\end{bmatrix}
=
\begin{bmatrix}
0 & 1 \\[4pt]
\omega^{2} & 0
\end{bmatrix}
\begin{bmatrix}
y(t) \\[4pt]
\dot{y}(t)
\end{bmatrix}
+
\begin{bmatrix}
0 \\[4pt]
\dfrac{1}{m_c}
\end{bmatrix}
\begin{bmatrix}
F_y
\end{bmatrix}
\end{equation}


% \subsubsection{Considerações finais}
% A aproximação relativa será conduzida em R-bar (eixo radial, $x$), sob realimentação de força, enquanto que a orientação do \emph{chaser} é controlada de forma independente para atender restrições de captura.
% Na próxima seção, apresentam-se as equações cinemáticas (quaternions) e dinâmicas (Euler) empregadas na malha de atitude e na definição da orientação relativa $q_{\mathrm{rel}}=q_t^{-1}\otimes q_c$.